\section{n2. Exact envelope transfer and every-degree approximation}
\textit{Supports: ENTIRE qualified G331; own n1.}

Use exactly n1's normal-target finite-domain inclusion, all source centres,
all supports and extended-primary gauges. No group action, flatness,
Galois transitivity, residue perfection, CM or Q-Gorenstein hypothesis is
imposed. For every finite continuous signed obstacle $f$ on $X$, put
$F=\rho^*f$. The raw geometric suprema satisfy the NEW identity
\[
                 \rho^*p_R(f)=p_S(F).                      \tag{5}
\]
This identity holds even if neither raw supremum is an envelope member.
In particular the source raw supremum of a pulled obstacle is constant
on EVERY full fibre. Continuity is not assumed to obtain this constancy.

We first reconstruct the canonical facts needed for the word envelope;
these are foundations, not novelty. On either nonempty compact section
set $B=\max(0,-\min f)$ and
\[
 A_m(f)=\{b:v(b)+mf(v)\ge0\text{ at EVERY point}\},\quad
 A_0=R,\quad h_m=-v(A_m)/m.
\]
For semilocal S the same tests use all its marked sections. These are
ideals containing $\mathfrak a^{\lceil mB\rceil}$, and
$A_mA_l\subset A_{m+l}$. They are unit-or-primary on R and boundary
ideals on S, with unit components allowed. A monic ideal-integrality
equation for $b$ violating one finite threshold would have its leading
term uniquely least-valued: its degree-i coefficient has value at least
$i(-mf(v))$. Thus the ideals are integrally closed, including all
infinite-value supports. Every $h_m$ is finite continuous, below $f$,
and actually negative effective ample on its own primary blowup BEFORE
limits. The unit case is zero.

Rational-profile density in G331 gives
\[
 p_C(f)=p_D(f)=\sup_m h_m,
 \qquad -B\le p_C(f)\le\min(f,0).                          \tag{6}
\]
For a C competitor choose a close rational profile and subtract a small
negative rational boundary constant; powers and a common denominator
give a rational profile below $f$. If $-v(I)/r\le f$ then
$I\subset A_r$. Arbitrarily small errors establish the C equality.
For a D competitor $g=\inf_i g_i\le f$, the increasing directed open
sets $\{g_i<f+\epsilon\}$ cover the compact section. A finite subcover
and directedness choose one index everywhere, and $g_i-\epsilon\in C$
is below $f$ and at least $g-\epsilon$. This proves the D equality,
including partial bottom. Constants $-B$ are boundary-ideal profiles.

Let $M_1=\|v(A_1)\|_\infty$ and $H_N=\max_{r\le N}h_r$.
For $m=kr+s$, $0\le s<r$, grading and nonpositivity imply
\[
 h_m\ge (kr/m)h_r-sM_1/m\ge h_r-rM_1/m,
 \qquad H_N-NM_1/m\le h_m\le p_C(f).
\]
Therefore EVERY degree converges pointwise to the finite LSC raw
supremum. If that supremum is continuous, compactness of the strict-test
cover $\{H_N>p_C-\epsilon\}$ yields uniform convergence of every $h_m$.
Conversely uniform convergence places it in C. A greatest D minorant
equals the raw supremum, so is USC as a decreasing completion, hence
continuous as well. Greatest C existence, greatest D existence, raw
continuity and uniform canonical convergence are consequently equivalent.
Negative constant shifts of competitors prove the obstacle map is
1-Lipschitz. No finite-generation premise or counterexample from
non-finite-generation is used.

To prove (5), pulling a target competitor up gives one inequality.
For every source C competitor $u\le F$, n1 gives an ACTUAL geometric
target competitor $\mathcal M_\rho u\in C_R$ with
$\mathcal M_\rho u\le f$ and $u\le\rho^*\mathcal M_\rho u$.
Taking suprema gives the reverse inequality. Equation (6) makes it
equally an identity for the raw D suprema. This proof uses the newly
proved positive maximum operator, not a declaration that a fibre extremum
is psh. It also proves the greatest-envelope equivalence for this fixed
obstacle on target and on its pullback source: (5), the compact quotient,
and (6) show continuity on one side iff on the other.

In particular suppose S has continuous geometric envelopes for every
finite continuous obstacle on its full boundary. Then R has the same
property, and
\[
      P_R f\in C_R,\qquad \rho^*P_R f=P_S(\rho^*f).          \tag{7}
\]
It dominates ALL target D competitors. Its canonical primary profiles
converge uniformly in EVERY degree and have actual curve-degree
positivity before convergence. The hypothesis can be weakened to source
continuity just for this pulled obstacle. This is a complete conditional
finite-cover theorem, not an assertion that every normal germ has such
a cover or that every regular source in positive/mixed characteristic
has the required continuity.

There is exact canonical contraction, with no extension equality:
\[
           A_m^R(f)=A_m^S(F)\cap R.                        \tag{8}
\]
Indeed every source value on an R element is its restricted value, and
restriction is surjective on the full sections. In completion the same
identity transports through G331's exact primary contraction and
boundary homeomorphisms, without a claim of unchanged generic degree or
of coefficient ideals commuting with every base change.

We prove a useful every-degree quantitative transfer. Let
$J_m=A_m^S(F)$ and $H=P_SF=\rho^*P_Rf$. Suppose a finite $K_0\ge0$
has actually been established such that
\[
        0\le H-h_m^S\le K_0/m \quad\text{on ALL of }Y.      \tag{9}
\]
Set $\mathcal N(J_m)=I_n(J_m)$, the last coefficient ideal in n1.
This is unit-or-primary without truncation: add generators $x$ of a
boundary power inside $J_m$ as redundant module generators; the pure
coefficient $x^n$ belongs to $I_n$. Its radical is therefore maximal
in the proper case. Product Gauss multiplicativity, on the SUPPORT
QUOTIENT field used in n1, gives
\[
 v(\mathcal N(J_m))=q\sum_{\sigma}w_\sigma(J_m),\qquad
 \kappa_m(v):=-v(\mathcal N(J_m))/(nm)
             ={1\over n}\int h_m^S\,d\mu_v.
\]
This is a primary profile with actual ample positivity. EVERY summand
is at most $f(v)$, so every coefficient obeys the canonical test and
$\mathcal N(J_m)\subset A_{nm}^R(f)$. Every summand is at least
$P_Rf(v)-K_0/m$, by (9) and (7). Hence
\[
        h_{nm}^R\ge\kappa_m\ge P_Rf-K_0/m.                 \tag{10}
\]
Generic inseparable lengths give the mass $n$; no ramification/residue
degree equality, flat determinant norm or equal source profiles on each
fibre is assumed. The profiles $h_m^S$ themselves need not be fibre
constant.

For $l=nm+s$, $0\le s<n$, $m\ge1$, target grading gives
\[
 h_l^R\ge (nm/l)h_{nm}^R-sM_1/l
       \ge P_Rf-(nK_0+sM_1)/l.
\]
The missing fraction of $P_Rf\le0$ only improves this bound.
For $1\le l<n$, $A_1^l\subset A_l$ gives $h_l\ge-M_1$.
Thus the single bound, on EVERY support and in EVERY degree, is
\[
           0\le P_Rf-h_l^R\le {n(K_0+M_1)\over l}.         \tag{11}
\]
For $f\ge0$ all canonical ideals are units and errors are zero.
For arbitrary continuous obstacles we claim uniform convergence, not
the existence of (9), a rational envelope, finite generation, or a rate.

There is also an exact maximum-based inclusion. Take
$N\ge\max(1,\lceil mB\rceil)$ in n1. Each
$I_l(J_m)+\mathfrak a^{lN}$ obeys the degree-$lm$ canonical tests,
by the root-value lower bound. Consequently
\[
 B_{J_m,N}\subset A_{D_nm}^R(f),\qquad
 h_{D_nm}^R\ge\mathcal M_\rho h_m^S.
\]
The norm construction is what keeps (11) at generic degree $n$.

For a decreasing net of finite continuous obstacles $f_i$ on a target
with this continuity property, $\inf_iP_Rf_i$ is an ACTUAL D member
and the greatest D minorant of $\inf_i f_i$. It is below the limit;
every competitor below that limit is below every $P_Rf_i$. Equation
(7) commutes with these decreasing nets and includes the bottom limit.
Gauge changes divide obstacles, profiles, raw and attained envelopes
by the same positive fibre denominator; the homogeneous cocycle is
G331's exact one. Dimension zero is the unique empty envelope. Products,
energy and unrestricted finite operators are separate targets.

Artifacts: notes/n2.tex and P12.3-normal-finite-cover-transfer.tex.
Rerun: python exact\_checks.py; python assemble\_delivery.py;
python verify\_delivery.py. The 9600 finite arithmetic checks of (11)
do not establish its geometric or source hypotheses.
